\documentclass[letterpaper, 12pt, parskip=full, DIV=10]{scrartcl}

\title{SOSC 13300: Social Science Inquiry III}
\subtitle{University of Chicago, Spring Quarter 2026.}

\RequirePackage{assets/template_MOW}
\RequirePackage{longtable}

\begin{document}
% The shared template loads bibentry, so keep a stub bibliography available.
\nobibliography{assets/bib}
\maketitle

\vskip -10 ex

\textbf{Location:} Pick Hall 222\\
\textbf{Course time:} Mon/Wed 3:00--4:20pm\\
\textbf{Course dates:} March 23--June 6, 2026\\
\textbf{Course GitHub:} \url{https://github.com/UChicago-pol-methods/SOSC13300-W26}

\textbf{Instructor:} Molly Offer-Westort; \href{mailto:mollyow@uchicago.edu}{mollyow@uchicago.edu}\\
\textbf{Office hours:} Book at \url{https://calendar.app.google/DZYbbZHQqCSYAuUr5}\\
\textbf{Office:} Pick Hall 328

\paragraph{Overview.}
This course is modified from a syllabus developed by Andy Eggers, which in turn built on a course created by Brendan Nyhan at Dartmouth College. In the final quarter of the Social Science Inquiry sequence, you will apply the research skills you developed in the first two quarters by contributing to a collaborative research project. The primary objective is to help you learn how to do research and how to critique research more effectively. A possible additional benefit is that our results may be strong enough to develop into a paper for submission, with students listed as coauthor---unless you decline that opportunity.

\paragraph{Course description.}
Our substantive focus is on climate change policy and politics, with special attention to how policy arguments are framed. Some arguments emphasize scientific information, some moral reasoning, and some economic consequences. Existing studies make different claims about which frames are effective, for whom they are effective, and whether some messages can even backfire. Early in the quarter, we will build enough substantive background to evaluate recent survey-based papers that we might replicate and extend. We will then choose a target paper, develop a research plan, pilot and refine our design, pre-register the study, field the survey, analyze the results, and write the paper together.

Each student will contribute to one part of the final study, and every student will write a version of the abstract, introduction, and discussion. At the end of the quarter, each student will also write a short reflection and critique of the combined draft. The course is designed to give you direct experience with the pleasures and frustrations of the research process while sharpening your ability to evaluate social science evidence.

\paragraph{Course objectives.}
This course should help you develop the ability to
\begin{itemize}
\item apply the tools you have learned so far in the sequence, including data analysis, programming, research design, and hypothesis testing;
\item use those tools and additional writing skills to produce a polished piece of research;
\item design and conduct an original experiment and analyze the results;
\item understand the differences between observational and experimental research; and
\item discuss important issues in climate change policy and politics.
\end{itemize}

\paragraph{Statistical software and computing.}
This version of the course repository is built around \texttt{R}, and you should expect to use \texttt{R} and RStudio throughout the quarter. You will need regular access to a computer with both installed. Some class meetings will include in-class computing work, so you should bring a laptop when possible. If access to a computer is a concern, please contact me and we will find a solution.

\paragraph{Assignments and grading.}
We will have a range of assignments in this course. Late work will not be accepted without prior permission.

\begin{center}
\begin{tabular}{l l}
\textbf{Component} & \textbf{Weight} \\
\hline
Class participation & 14\% \\
Research proposal & 15\% \\
Section write-up (theory, methods, or results) & 15\% \\
Other assignments & 8\% each, 56\% total \\
\end{tabular}
\end{center}

\paragraph{Assignment files.}
Full assignment prompts live in the course repository's assignments folder. 

\paragraph{Accommodations.}
Please reach out if you would like to request accommodations that would better facilitate your learning in this course. Student Disability Services is also available to provide resources and support and may approve specific academic accommodations: \url{https://disabilities.uchicago.edu/}.

\paragraph{Academic honesty and AI use.}
All students are expected to be familiar with and follow the University's academic honesty policies: \url{https://studentmanual.uchicago.edu/academic-policies/academic-honesty-plagiarism/}. 
For this course, you may use generative artificial intelligence tools such as Claude or Codex for generating code. 
No citation is required for code generated by AI, but you are responsible for understanding the code and verifying that it works correctly. 
You may also use AI tools for brainstorming, but you may not directly use AI-generated prose in course assignments. 
Because AI systems can hallucinate information, you should verify any factual claims they provide against reliable sources.

\section*{Course outline}

\small
\setlength{\LTpre}{0pt}
\setlength{\LTpost}{0pt}
\renewcommand{\arraystretch}{1.1}
\begin{longtable}{@{}>{\raggedright\arraybackslash}p{0.27\textwidth}>{\raggedright\arraybackslash}p{0.27\textwidth}>{\raggedright\arraybackslash}p{0.38\textwidth}@{}}
\textbf{Monday class} & \textbf{Wednesday class} & \textbf{Assignment} \\
\midrule
\endhead
(3/23) Introduction; science and social science of climate change &
(3/25) Science and social science of climate policy &
(3/27) Problem set 1: data analysis refresher \\

(3/30) Discussion of replication/extension candidates, part 1 &
(4/1) Qualtrics for survey development &
(4/3) Problem set 2: Qualtrics exercise \\

(4/6) Discussion of replication/extension candidates, part 2 &
No class &
(4/10) Problem set 3: data analysis from Qualtrics output \\

(4/13) Research design and adaptive experimentation &
(4/15) Discuss research plan; allocation of responsibilities &
Due Tue 4/14 at 3pm: proposal for our class research project (worth 15\%) \\

(4/20) Group work for pilot &
(4/22) Group work for pilot; data visualization &
(4/24) Draft of group output (survey, literature spreadsheet, or PAP components). Pilot fielded over the weekend. \\

(4/27) Group work for PAP &
(4/29) Group work for PAP &
(5/1) Final version of group output using pilot results. PAP submitted and survey fielded over the weekend. \\

(5/4) Academic writing: theory, methods, and results &
(5/6) Optional meeting for writing consult &
(5/8) Section draft (theory, methods, or results), worth 15\%. Partial combined draft distributed by Mon 5/11 at 5pm. \\

(5/11) Academic writing: introduction, abstract, and discussion &
(5/13) Optional meeting for writing consult &
(5/15) Remainder draft (abstract, introduction, and discussion). Full combined draft distributed by Wed 5/20 at 5pm. \\

(5/18) Presenting academic research &
(5/20) Presentation of results &
(5/29) Reflection and critique due \\
\end{longtable}
\normalsize

\section*{Detailed schedule}

\subsection*{Class 1.1 Monday, March 23: Course orientation; the science and social science of climate change}
\begin{itemize}
\item Overview of course objectives and structure: producing a publishable research paper on the politics of climate change.
\end{itemize}

\textbf{Readings.}
\begin{itemize}
\item None.
\end{itemize}


\subsection*{Class 1.2 Wednesday, March 25: The science and social science of framing}
\begin{itemize}
\item An introduction to social framing and its possible role in advancing climate change policy.
\end{itemize}

\textbf{Readings.}
\begin{itemize}
\item \bibentry{sparkman2022americans}
\item \bibentry{chiancone2024goodnews}
\item \bibentry{climatereality2024ignore}
\end{itemize}

\textbf{Assignment due Friday, March 27 at 3pm on Canvas.}
See \path{assignment01-data-analysis-refresher.R} for the full prompt.

\subsection*{Class 2.1 Monday, March 30: Discussion of replication/extension candidates, part 1}
\begin{itemize}
\item Closely examine papers that we could replicate and extend in our own research. Part I focuses on background and framing the problem.
\end{itemize}

\textbf{Readings.}
\begin{itemize}
\item \bibentry{swirethompson2020backfire}
\item \bibentry{druckman2019motivated}
\item \bibentry{nyhan2021durability}
\end{itemize}

\subsection*{Class 2.2 Wednesday, April 1: Qualtrics for survey development}
\begin{itemize}
\item Develop facility with Qualtrics as a tool for collecting data for a survey experiment.
\end{itemize}

\textbf{Readings.}
\begin{itemize}
\item \bibentry{diaz2020survey}
\item \bibentry{leeper2018mistake}
\end{itemize}

\textbf{Assignment due Friday, April 3 at 3pm via Canvas.}
See \path{assignment02-qualtrics-exercise.md} for the full prompt.

\subsection*{Class 3.1 Monday, April 6: Discussion of replication/extension candidates, part 2}
Before class, be ready to
\begin{itemize}
\item explain in simple terms what each study did and found;
\item highlight something confusing, unclear, or unsatisfying about one of the papers; and
\item assess whether it would be useful to replicate or extend these papers.
\end{itemize}

\textbf{Readings.}
\begin{itemize}
\item \bibentry{goldwert2026megastudy}
\item \bibentry{vlasceanu2024addressing}
\item \bibentry{voelkel2026registered}
\end{itemize}

\subsection*{Class 3.2 Wednesday, April 8: Project scoping and proposal workshop}

\textbf{Readings.}
\begin{itemize}
\item None.
\end{itemize}

\textbf{Assignment due Friday, April 10 at 3pm via Canvas.}
See \path{assignment03-qualtrics-output-analysis.R} for the full prompt.

\subsection*{Class 4.1 Monday, April 13: Research design and adaptive experimentation}
\textbf{Readings.}
\begin{itemize}
\item \bibentry{gerber2012field} Chapter 13: ``Writing an Experimental Proposal, Research Report, and Journal Article.''
\end{itemize}

\textbf{Assignment due Tuesday, April 14 at 3pm via Canvas.}
See \path{assignment04-research-proposal.md} for the full prompt and team descriptions.

\subsection*{Class 4.2 Wednesday, April 15: Starting work on our research}
\begin{itemize}
\item Discuss the decision about research focus and our overall plan.
\item Begin work in groups, with short group reports at the end of class.
\item Review group assignments, deliverables, and the research design.
\end{itemize}

\textbf{Readings.}
\begin{itemize}
\item \bibentry{monogan2015preregistration}
\item \bibentry{chen2021preanalysis}
\end{itemize}

\textbf{Additional reading.}
\begin{itemize}
\item \bibentry{banerjee2020moderation}
\end{itemize}

\subsection*{Class 5.1 Monday, April 20: Group work on draft PAP components}
\begin{itemize}
\item Continue group work, with short group reports at the end of class.
\end{itemize}

\subsection*{Class 5.2 Wednesday, April 22: Group work on draft PAP components}
\begin{itemize}
\item Continue group work, with short group reports at the end of class.
\end{itemize}

\textbf{Assignment due Friday, April 24 at 3pm via Google Docs.}
See \path{assignment05-group-draft-pap-components.md} for the full prompt.

\subsection*{Class 6.1 Monday, April 27: Group work on PAP components}
\begin{itemize}
\item Work with pilot data, with short group reports at the end of class.
\end{itemize}

\subsection*{Class 6.2 Wednesday, April 29: Group work on PAP components}
\begin{itemize}
\item Continue group work, with short group reports at the end of class.
\end{itemize}

\textbf{Assignment due Friday, May 1 at 3pm via Google Docs.}
See \path{assignment06-group-final-pap-components.md} for the full prompt.

\subsection*{Class 7.1 Monday, May 4: Academic writing: theory, methods, and results}
\textbf{Readings for reference.}
\begin{itemize}
\item \bibentry{mensh2017structuring}
\item \bibentry{coppock2018writeup}
\item \bibentry{coppock2022digital}
\end{itemize}

\textbf{Assignment due Friday, May 8 at 3pm via Canvas.}
See \path{assignment07-section-draft.md} for the full prompt.

\subsection*{Class 7.2 Wednesday, May 6: Optional meeting for writing consult}
Attendance is not required, with the expectation that your group will meet outside of class time. Your group can also use class as a working session.

\subsection*{Class 8.1 Monday, May 11: Academic writing: abstract, introduction, and discussion}
\textbf{Readings.}
\begin{itemize}
\item \bibentry{chockalingam2021limited}
\end{itemize}

The partial combined draft will be distributed by Monday, May 11 at 5pm.

\subsection*{Class 8.2 Wednesday, May 13: Optional meeting for writing consult}
Attendance is not required, with the expectation that your group will meet outside of class time. Your group can also use class as a working session.

\textbf{Assignment due Friday, May 15 at 3pm via Canvas.}
See \path{assignment08-abstract-introduction-discussion.md} for the full prompt.

\subsection*{Class 9.1 Monday, May 18: Presenting research}
\begin{itemize}
\item Data visualization and presentation of research findings.
\end{itemize}

\subsection*{Class 9.2 Wednesday, May 20: Discussion of our research}
The full combined draft will be distributed by Wednesday, May 20 at 5pm.

\textbf{Assignment due Friday, May 29 at 3pm via Canvas.}
See \path{assignment09-reflection-and-critique.md} for the full prompt.

\end{document}
